\title{DeepSeekMath: Pushing the Limits of Mathematical Reasoning in Open Language Models}

\begin{abstract}
Mathematical reasoning remains a substantial obstacle for language models owing to its inherently intricate and structured demands. This study presents DeepSeekMath~7B, which is developed by continuing the pre-training of DeepSeek-Coder-Base-v1.5 7B on a curated dataset comprising 120B math-centric tokens aggregated from Common Crawl, in addition to accompanying natural language and code resources. DeepSeekMath~7B attains a remarkable 51.7\% on the competitive MATH benchmark, without the aid of supplementary toolkits or ensemble voting, reaching performance levels close to those of Gemini-Ultra and GPT-4. Evaluating self-consistency with 64 sampled outputs from DeepSeekMath~7B yields a 60.9\% score on MATH. The notable proficiency in mathematical reasoning demonstrated by DeepSeekMath~7B arises from two principal innovations: firstly, the deployment of a rigorously designed data filtration pipeline that exploits publicly accessible web sources; secondly, the introduction of Group Relative Policy Optimization (GRPO)—a modified approach to Proximal Policy Optimization (PPO)—which strengthens mathematical reasoning capacities while improving PPO's memory efficiency.
\end{abstract}

\section{Introduction}

Large language models (LLM) have transformed the landscape of mathematical reasoning within artificial intelligence, driving notable progress in benchmarks for quantitative reasoning \citep{MATH} and geometric reasoning \citep{trinh2024solving}. In addition, these models have demonstrated considerable utility in aiding humans to tackle intricate mathematical challenges \citep{tao}. Nevertheless, state-of-the-art systems like GPT-4 \citep{gpt4} and Gemini-Ultra \citep{gemini} remain unavailable to the public, and existing open-source alternatives exhibit significantly lower levels of performance.

In this work, we present DeepSeekMath, a specialized language model designed to substantially surpass the mathematical reasoning abilities of other open-source models and closely approach the academic benchmark performance of GPT-4. Central to this achievement is the development of the DeepSeekMath Corpus—a large-scale, high-quality pre-training dataset containing 120B math tokens. This corpus is constructed from Common Crawl (CC) data through the use of a fastText-based classifier \citep{joulin2016fasttext}. Initially, the classifier is trained on positive samples drawn from OpenWebMath \citep{openwebmath}, complemented by a wide array of negative examples sampled from general web pages. The trained classifier is then deployed to identify additional positive samples from CC, after which these instances undergo further filtering via human annotation. The resulting improved dataset is used to further train and enhance the classifier. Evaluation demonstrates the high quality of the resulting corpus, as evidenced by the base model DeepSeekMath-Base 7B achieving a 64.2\% score on GSM8K \citep{gsm8k} and 36.2\% on the MATH dataset \citep{MATH}, surpassing the performance of Minerva 540B \citep{minerva}. Furthermore, the multilingual nature of the DeepSeekMath Corpus enables enhanced outcomes on Chinese mathematical benchmarks \citep{wei2023cmath,agieval}. We consider our methodologies in mathematical data construction a valuable foundation for the research field, acknowledging that there remains considerable scope for further advancement.

DeepSeekMath-Base is constructed by initializing with DeepSeek-Coder-Base-v1.5 7B \citep{deepseek-coder}, as we find that leveraging a model pre-trained on code yields superior results to beginning with a general-purpose LLM. Additionally, our findings indicate that mathematical pre-training not only strengthens performance on mathematical tasks but also brings improvements on broader benchmarks such as MMLU \citep{mmlu} and BBH \citep{bbh}. This suggests a positive transfer from mathematical training to general reasoning abilities.

Following the pre-training phase, we subject DeepSeekMath-Base to mathematical instruction tuning, incorporating data from chain-of-thought \citep{cot}, program-of-thought \citep{pot,pal}, and tool-integrated reasoning \citep{tora} approaches. The finalized model, DeepSeekMath-Instruct 7B, surpasses all other 7B models and demonstrates performance on par with 70B open-source instruction-tuned models.

In addition, we propose Group Relative Policy Optimization (GRPO), a reinforcement learning algorithm inspired by Proximal Policy Optimization (PPO) \citep{schulman2017proximal}, but with key differences. Unlike PPO, GRPO eliminates the need for a critic network, relying instead on baselines computed from group performance, which greatly cuts down on the computational resources required for training. Leveraging only a limited selection of English instruction tuning data, GRPO achieves marked performance gains relative to the already strong DeepSeekMath-Instruct baseline, with improvements observed in both in-domain tasks (GSM8K: 82.9\% $\rightarrow$ 88.2\%, MATH: 46.8\% $\rightarrow$ 51.7\%) and generalization to out-of-domain mathematics benchmarks (e.g., CMATH: 84.6\% $\rightarrow$ 88.8\%) during reinforcement learning. We also present a unified perspective that encompasses methods like Rejection Sampling Fine-Tuning (RFT) \citep{yuan2023scaling}, Direct Preference Optimization (DPO) \citep{dpo}, PPO, and our GRPO, conceptualizing them all within the spectrum of direct or streamlined reinforcement learning strategies. To examine the core features underlying this paradigm, we carry out thorough experimentation, such as comparing online and offline learning, outcome-oriented and process-oriented supervision, as well as single-step and iterative RL frameworks. Finally, we discuss the underlying reasons that reinforcement learning enhances the capabilities of instruction-tuned models and outline prospective directions for advancing RL effectiveness grounded in this unified framework.

\subsection{Contributions}
Our contribution includes scalable math pre-training, along with the exploration and analysis of reinforcement learning.

\textbf{Math Pre-Training at Scale}

\begin{itemize}[topsep=0pt]
    \item We demonstrate that the openly available Common Crawl dataset harbors significant resources relevant to mathematics. Through the application of a carefully curated data filtering process, we assembled the DeepSeekMath Corpus—a large-scale, high-quality mathematical dataset consisting of 120B tokens extracted from web pages selected for their mathematical content. This dataset is nearly sevenfold larger than the math web pages employed by Minerva \citep{minerva}, and exceeds the size of the recently introduced OpenWebMath \citep{openwebmath} by a factor of nine.

\item DeepSeekMath-Base 7B, our foundational pre-trained model, demonstrates performance on par with Minerva 540B \citep{minerva}, suggesting that model size alone does not determine proficiency in mathematical reasoning. Robust results can also be attained by smaller models when they are trained on curated, high-quality datasets.

\item We present results from our experiments involving math training. Pretraining models on code enhances their performance on mathematical tasks, regardless of whether they employ external tools. This provides a partial response to the ongoing debate: \textit{does code training enhance reasoning skills?} Our findings indicate that it does, particularly in the realm of mathematical reasoning.

    \item Although training on arXiv papers is common, especially in many math-related papers, it brings no notable improvements on all mathematical benchmarks adopted in this paper.
\end{itemize}

\textbf{Exploration and Analysis of Reinforcement Learning}

\begin{itemize}[topsep=0pt]
    \item We present Group Relative Policy Optimization (GRPO), a novel reinforcement learning algorithm that is both resource-efficient and effective. By omitting the critic network and deriving the baseline directly from group score information, GRPO greatly reduces the computational overhead typically associated with Proximal Policy Optimization (PPO).
    \item Our experiments demonstrate that applying GRPO leads to marked improvements in the instruction-tuned model DeepSeekMath-Instruct when relying solely on instruction-tuning datasets. In addition, we note gains in generalization, with superior out-of-domain performance observed during reinforcement learning.
    \item We establish a unified framework for interpreting a variety of methodologies, including RFT, DPO, PPO, and GRPO. Through comprehensive empirical studies—such as comparing online versus offline training, evaluating outcome against process supervision, and analyzing both single-turn and iterative reinforcement learning—we scrutinize key components underlying this paradigm.
    \item Leveraging this unified framework, we analyze the mechanisms that contribute to the efficacy of reinforcement learning and highlight several prospective pathways for developing more robust reinforcement learning strategies for LLMs.
\end{itemize}

\subsection{Summary of Evaluations and Metrics}

\begin{itemize}[topsep=0pt]
    \item \textbf{English and Chinese Mathematical Reasoning}: Our evaluation framework encompasses a wide range of English and Chinese benchmarks, featuring mathematical questions spanning from elementary to collegiate levels. The English evaluation suite consists of GSM8K \citep{gsm8k}, MATH \citep{MATH}, SAT \citep{llemma}, OCW Courses \citep{minerva}, and MMLU-STEM \citep{mmlu}. For Chinese, we employ MGSM-zh \citep{mgsm}, CMATH \citep{wei2023cmath}, Gaokao-MathCloze \citep{agieval}, and Gaokao-MathQA \citep{agieval}. The evaluation focuses on two core capabilities: the production of fully self-contained written solutions without external tool assistance, as well as successful problem solving through Python.

Across English evaluation datasets, DeepSeekMath-Base achieves results on par with the proprietary Minerva 540B \citep{minerva}, and outperforms all publicly available base models such as Mistral 7B \citep{mistral} and Llemma-34B \citep{llemma}, whether or not these have received specialized mathematical pre-training—and frequently does so by a notable margin. Importantly, DeepSeekMath-Base attains even greater results on Chinese test sets, a likely consequence of including high-quality non-English data in our pre-training, in contrast with prior research \citep{minerva,llemma} that limits data sources to English-only math content. Subsequent instruction tuning for mathematics and reinforcement learning yield DeepSeekMath-Instruct and DeepSeekMath-RL, both of which display robust capabilities—achieving, for the first time in the open-source domain, accuracy above 50\% on the challenging, competition-level MATH benchmark.

\item \textbf{Formal Mathematics}: We assess the capabilities of DeepSeekMath-Base on the informal-to-formal theorem proving task, as outlined in \citep{dsp_proof}, utilizing the miniF2F benchmark \citep{minif2f} with Isabelle \citep{isabelle} designated as the proof assistant. In these evaluations, DeepSeekMath-Base achieves robust performance in few-shot autoformalization scenarios.
\item \textbf{Natural Language Understanding, Reasoning, and Code}: To provide a holistic assessment of the models’ proficiency in comprehension, reasoning, and coding, we subject DeepSeekMath-Base to a variety of benchmarks: the Massive Multitask Language Understanding (MMLU) evaluation \citep{mmlu} featuring 57 multiple-choice tests spanning a wide range of topics, BIG-Bench Hard (BBH) \citep{bbh} comprising 23 tasks emphasizing multi-step reasoning, and both HumanEval \citep{codex} and MBPP \citep{mbpp} for code-related performance measurement. The pre-training focused on mathematics is observed to enhance the model’s effectiveness in both language understanding and complex reasoning.

\section{Math Pre-Training}

\subsection{Data Collection and Decontamination} 

Here, we describe the methodology used to build the DeepSeekMath Corpus using data sourced from Common Crawl. Figure \ref{fig:data_collect} illustrates the iterative workflow we employ to methodically assemble an extensive mathematical dataset from Common Crawl, beginning with an initial seed corpus—typically a limited yet curated set of math-focused data. Importantly, this iterative framework is adaptable and can be effectively utilized for domains beyond mathematics, such as software code.

Our approach begins with OpenWebMath \citep{openwebmath}—a curated set of high-quality mathematical texts from the web—as the foundational seed corpus. Utilizing this dataset, we construct a fastText model \citep{joulin2016fasttext} to retrieve additional web pages resembling those in OpenWebMath. For model training, we randomly sample 500,000 examples from the seed corpus as positive instances and another 500,000 pages from Common Crawl as negative samples. The training process is conducted using an open-source fastText implementation\footnote{\url{https://fasttext.cc}}, with parameters set as follows: vector dimension of 256, learning rate of 0.1, a maximum word n-gram length of 3, a minimum frequency threshold of 3 for words, and three epochs. To shrink the vast Common Crawl data, we use both exact URL-based and near-deduplication approaches, ultimately curating a set of 40 billion HTML web pages. We subsequently apply the trained fastText model to this deduplicated collection, retrieving mathematical pages according to their model-predicted scores and retaining only the highest-scoring content to ensure quality. The volume of retained data is validated through pre-training tests spanning the top 40B, 80B, 120B, and 160B tokens. During the initial cycle, we opt to maintain the top 40B tokens.

Following the initial round of data collection, a substantial number of mathematical web pages remain uncollected. This is primarily due to the limited diversity within the positive samples used to train the fastText model. To remedy this, we seek out and incorporate additional mathematics-focused web sources to further diversify the seed set, with the aim of refining the fastText classifier. To this end, we segment the full Common Crawl dataset by distinct domains—each domain comprising web pages with the same base URL. We then determine, for every domain, what fraction of its web pages were obtained during the first collection phase. Domains in which more than 10\% of the pages were retrieved are designated as math-focused (for example, \url{mathoverflow.net}). For such domains, we proceed to manually label the specific URLs corresponding to mathematical content (e.g., \url{mathoverflow.net/questions}). Any yet-uncollected web pages found under these annotated URLs are added to the seed set. This strategy results in an expanded collection of positive training samples, allowing the fastText model to become more effective at retrieving mathematical pages in subsequent iterations. After running this iterative data acquisition process four times, we accumulate a total of 35.5M mathematical web pages, amounting to 120B tokens. In the fourth cycle, we observe that close to 98\% of the available content had already been gathered by the third iteration, leading us to discontinue the data collection process.

In order to prevent benchmark contamination, we adopt the approach described in \cite{deepseek-coder}, eliminating web pages that include either questions or answers from established English mathematical benchmarks like GSM8K \citep{gsm8k} and MATH \citep{MATH}, as well as Chinese benchmarks such as CMATH \citep{wei2023cmath} and AGIEval \citep{agieval}. Specifically, our filtering strategy involves removing any segment of text from the mathematical training set that contains a 10-gram sequence perfectly matching any substring present in the evaluation benchmarks. For benchmark texts between 3 and 9 grams in length, we also perform an exact string match to screen out any web pages that might introduce contamination.

\subsection{Validating the Quality of the DeepSeekMath~Corpus}

We run pre-training experiments to investigate how the DeepSeekMath~Corpus is compared with the recently released math-training corpora:

\begin{itemize}[topsep=0pt]
    \item \textbf{MathPile} \citep{mathpile}: This dataset comprises 8.9B tokens collected from diverse sources such as textbooks, Wikipedia, ProofWiki, CommonCrawl, StackExchange, and arXiv, with arXiv comprising more than 85\% of the overall content;
    \item \textbf{OpenWebMath} \citep{openwebmath}: Derived from CommonCrawl, this 13.6B-token corpus has been filtered to focus exclusively on mathematics-related material;
    \item \textbf{Proof-Pile-2} \citep{llemma}: Incorporating OpenWebMath, AlgebraicStack (which contributes 10.3B tokens from mathematical code), and arXiv (accounting for 28.0B tokens), this large mathematical dataset follows the source proportion of 2:4:1 (arXiv:Web:Code) outlined in \cite{llemma} when utilized for experiments.
\end{itemize}

\subsubsection{Training Setting}
\label{sec:quality-policy}
Mathematics-focused fine-tuning is performed on a base language model comprising 1.3B parameters, architecturally aligned with the DeepSeek LLM family \citep{deepseek-llm} and identified as DeepSeek-LLM 1.3B. Each mathematical dataset is used to individually train a model for a total of 150B tokens. The entirety of training takes advantage of HAI-LLM \citep{haillm}, an optimized and lightweight framework. In accordance with established procedures for DeepSeek LLMs, the AdamW optimizer \citep{adamW} is employed, set with $\beta_1=0.9$, $\beta_2=0.95$, and $\mathrm{weight\_decay}=0.1$. The learning rate follows a multi-stage decay policy: it reaches its maximum after 2{,}000 warmup steps, drops to 31.6\% of the peak by 80\% of the total training, and reduces further to 10.0\% of the maximum value by the 90\% completion mark. The highest learning rate applied is 5.3e-4, using batches comprising 4 million tokens and a maximum context length of 4K.

\subsubsection{Evaluation Results}

\textbf{The DeepSeekMath~Corpus is of high quality, covers multilingual mathematical content, and is the largest in size.}

\begin{itemize}[topsep=0pt]
    \item \textbf{High-quality}:  
    Downstream task results on eight mathematical benchmarks were assessed using few-shot chain-of-thought prompting~\cite{cot}. As reported in Table~\ref{tab:corpora_comparison}, models trained with the DeepSeekMath~Corpus consistently outperform counterparts. Additionally, Figure~\ref{fig:corpora_comparisons} illustrates that, at the 50B token mark (equivalent to one full pass through Proof-Pile-2), the DeepSeekMath-trained model surpasses one trained on Proof-Pile-2, underscoring the superior average quality of the DeepSeekMath~Corpus. 

    \item \textbf{Multilingual}:  
    DeepSeekMath~Corpus offers diverse linguistic coverage, with English and Chinese comprising the majority of its data. Table~\ref{tab:corpora_comparison} reveals that, when models are trained on DeepSeekMath~Corpus, mathematical reasoning abilities in both English and Chinese improve significantly. In contrast, previously available mathematical datasets, which are heavily focused on English, produce only marginal benefits or may even detract from Chinese math reasoning proficiency.

    \item \textbf{Large-scale}:  
    In terms of data size, DeepSeekMath~Corpus greatly exceeds the volume of prior mathematical corpora. Figure~\ref{fig:corpora_comparisons} demonstrates that DeepSeek-LLM~1.3B, once trained on the DeepSeekMath~Corpus, enjoys a notably accelerated and sustained progression in learning. On the other hand, baseline datasets are much more limited in scope, have undergone multiple passes during model training, and exhibit early saturation in model improvement.
\end{itemize}

\subsection{Training and Evaluating DeepSeekMath-Base 7B}

This section presents DeepSeekMath-Base 7B, a foundational model that exhibits notable proficiency in reasoning tasks, particularly those involving mathematics. The model is built upon DeepSeek-Coder-Base-v1.5 7B \citep{deepseek-coder} as its initialization and undergoes training over 500B tokens. The data composition includes 56\% from the DeepSeekMath Corpus, 4\% sourced from AlgebraicStack, 10\% from arXiv, 20\% derived from Github code, with the final 10\% comprised of natural language data—sourced from Common Crawl in both English and Chinese. The training configuration primarily follows the guidelines described in Section \ref{sec:quality-policy}; however, we adjust the learning rate's maximum value to 4.2e-4 and employ a batch size encompassing 10M tokens.

We present an extensive evaluation of DeepSeekMath-Base 7B's mathematical proficiency, examining its competence in independently generating complete mathematical solutions, leveraging tools to address mathematical tasks, and performing rigorous formal theorem proving. Additionally, we outline a broader characterization of the base model, detailing its capabilities in natural language comprehension, reasoning, and programming beyond mathematical problem solving.

\paragraph{Mathematical Problem Solving with Step-by-Step Reasoning}

We assess the ability of DeepSeekMath-Base to solve mathematical problems through few-shot chain-of-thought prompting \citep{cot}, drawing on eight different benchmarks presented in both English and Chinese. These benchmarks feature tasks involving quantitative reasoning—such as those in GSM8K \citep{gsm8k}, MATH \citep{MATH}, and CMATH \citep{wei2023cmath}—as well as multiple-choice questions including MMLU-STEM \citep{mmlu} and Gaokao-MathQA \citep{agieval}. The selected datasets span a wide spectrum of mathematical areas, ranging from elementary topics to those encountered at the college level.

Table \ref{tab:base_math_cot} illustrates that among open-source base models, DeepSeekMath-Base 7B consistently achieves the highest performance across all eight evaluated benchmarks. This comparison includes prominent models such as the widely adopted general-purpose Mistral 7B \citep{mistral} and the recently introduced Llemma 34B \citep{llemma}, the latter of which was trained on Proof-Pile-2 \citep{llemma} for mathematical proficiency. Remarkably, DeepSeekMath-Base demonstrates more than a 10\% absolute advantage over its open-source peers on the challenging MATH dataset. Furthermore, it exceeds the performance of Minerva 540B \citep{minerva}—a closed-source base model 77 times its size, based on PaLM \citep{palm} and additionally trained on mathematics-related texts.

\paragraph{Mathematical Problem Solving with Tool Use} Our assessment of program-assisted mathematical reasoning involves the GSM8K and MATH benchmarks, employing few-shot program-of-thought prompting as outlined by \citep{pot,pal}. In this setup, models are asked to address each mathematical question by generating a Python script, with access to libraries such as \textit{math} and \textit{sympy} for complex computations. The answer is determined by running the program and checking its output. According to Table \ref{tab:base_math_pot_proof}, DeepSeekMath-Base 7B achieves superior performance over the previous leading model, Llemma 34B.

\paragraph{Formal Mathematics}

The automation of formal proof plays a crucial role in guaranteeing both the correctness and trustworthiness of mathematical proofs, while also streamlining the proving process. Interest in this area has surged in recent years. We assess DeepSeekMath-Base 7B using the informal-to-formal proof task presented by \citep{dsp_proof}, which involves generating a formalized proof from an informal proposition, its formal version, and an accompanying informal proof. Our evaluation utilizes the miniF2F benchmark \citep{minif2f}, which focuses on Olympiad-level mathematics. For each problem, we apply few-shot prompting to produce a formal proof in Isabelle. In accordance with the methodology outlined in \cite{dsp_proof}, our approach tasks the models with drafting proof outlines, and then employs the existing automated prover Sledgehammer \citep{sledgehammer} to complete the details. Table \ref{tab:base_math_pot_proof} highlights that DeepSeekMath-Base 7B achieves robust results in the automatic formalization of proofs.

\paragraph{Natural Language Understanding, Reasoning, and Code}

We assess the capabilities of the models in natural language understanding using MMLU \citep{mmlu}, in reasoning using BBH \citep{bbh}, and in programming through HumanEval \citep{codex} and MBPP \citep{mbpp}. As summarized in Table \ref{tab:understanding_reasoning_code}, DeepSeekMath-Base 7B achieves marked improvements over its predecessor, DeepSeek-Coder-Base-v1.5 \citep{deepseek-coder}, on MMLU and BBH, demonstrating the beneficial effects of mathematics-focused training on both language comprehension and reasoning abilities. Furthermore, due to the integration of code tokens during continual training, DeepSeekMath-Base 7B is able to preserve the coding performance level established by DeepSeek-Coder-Base-v1.5 on both programming benchmarks. In summary, DeepSeekMath-Base 7B substantially surpasses the general-purpose Mistral 7B model \citep{mistral} across all three benchmarks related to reasoning and coding.

\section{Supervised Fine-Tuning}

\subsection{SFT Data Curation}

A diverse instruction-tuning dataset for mathematics was developed, encompassing English and Chinese questions drawn from a range of mathematical disciplines and complexities. Each problem is accompanied by solutions formatted as chain-of-thought (CoT) \citep{cot}, program-of-thought (PoT) \citep{pot,pal}, and tool-integrated reasoning formats \citep{tora}. In total, the dataset contains 776,000 training instances.

\begin{itemize}[topsep=0pt]
    \item \textbf{English mathematical datasets}: We enrich GSM8K and MATH with annotations featuring tool-based solution traces and further include a curated portion of MathInstruct \citep{MathInstruct}, as well as the training portion of Lila-OOD \citep{lila}, both employing CoT or PoT solution methodologies. Our English dataset compilation spans a wide array of mathematical domains, including algebra, probability, number theory, calculus, and geometry.
    \item \textbf{Chinese mathematical datasets}: We aggregate Chinese K-12 mathematics questions across 76 specific subfields—for example, linear equations—and provide annotated answers using both CoT and tool-based reasoning approaches.
\end{itemize}

\subsection{Training and Evaluating DeepSeekMath-Instruct 7B}

This section presents DeepSeekMath-Instruct 7B, a model refined via mathematical instruction tuning using the DeepSeekMath-Base as its foundation. During training, example problems are combined in random order until a total sequence length of up to 4K tokens is reached. The model undergoes 500 training steps, employing a batch size of 256 and maintaining a fixed learning rate of 5e-5.

For evaluation, we assess the model’s mathematical reasoning abilities in both tool-assisted and tool-free scenarios, utilizing four quantitative reasoning benchmarks in both English and Chinese. Our assessment includes a comparison to the state-of-the-art models available at the time:

\begin{itemize}[topsep=0pt]
    \item \textbf{Closed-source models} consist of: (1) the GPT series, with GPT-4 \citep{gpt4} and GPT-4 Code Interpreter~\footnote{\url{https://openai.com/blog/chatgpt-plugins\#\#code-interpreter}} representing the leading-edge variants, (2) Gemini Ultra and Pro \citep{gemini}, (3) Inflection-2 \citep{inflection-2}, (4) Grok-1~\footnote{\url{https://x.ai/model-card}}, and a range of more recent offerings from Chinese organizations, including (5) Baichuan-3~\footnote{\url{https://www.baichuan-ai.com}} and (6) the newest version of GLM, GLM-4~\footnote{\url{https://open.bigmodel.cn/dev/api\#glm-4}}

    from the GLM family \citep{glm}.

These models are designed for broad, general-purpose applications and have typically been subjected to various alignment techniques.
    \item \textbf{Open-source models} comprise: general-purpose systems such as (1) DeepSeek-LLM-Chat 67B \citep{deepseek-llm}, (2) Qwen 72B \citep{qwen}, (3) SeaLLM-v2 7B \citep{seallm}, and (4) ChatGLM3 6B \citep{chatglm3}, along with models that specifically target mathematical competencies, including (5) InternLM2-Math 20B~\footnote{\url{https://github.com/InternLM/InternLM-Math}}, an extension of InternLM2 that incorporates both mathematics-specific training and instruction tuning, (6) Math-Shepherd-Mistral 7B, which applies proximal policy optimization (PPO) training \citep{schulman2017proximal} to Mistral 7B \citep{mistral} using a reward model guided by process supervision, (7) WizardMath models \citep{wizardmath}, which enhance Mistral 7B and Llama-2 70B \citep{llama2} in mathematical reasoning by employing the evolve-instruct approach—an instruction tuning method based on AI-evolved instructions—and PPO training, with the majority of training data derived from GSM8K and MATH benchmarks, (8) MetaMath 70B \citep{metamath}, which is created by fine-tuning Llama-2 70B on expanded versions of GSM8K and MATH datasets, (9) ToRA 34B \cite{tora}, a CodeLlama 34B model adjusted for tool-integrated mathematical reasoning, and (10) MAmmoTH 70B \citep{MathInstruct}, which results from instruction-tuning Llama-2 70B on the MathInstruct dataset.
\end{itemize}

As presented in Table \ref{tab:sft_rl_math}, when tool usage is prohibited, DeepSeekMath-Instruct 7B achieves notable results in step-by-step reasoning tasks. In particular, on the challenging MATH dataset, our model outperforms all publicly available models as well as most closed-source counterparts (such as Inflection-2 and Gemini Pro) by a margin of no less than 9\% in absolute terms. This advantage persists even when compared to models of considerably greater size (for example, Qwen 72B) or those that have undergone specialized reinforcement learning targeting mathematical abilities (such as WizardMath-v1.1 7B). DeepSeekMath-Instruct is competitive with Chinese proprietary models like GLM-4 and Baichuan-3 on MATH; however, its performance remains lower than that of leading models such as GPT-4 and Gemini Ultra.

In an evaluation framework that permits models to combine natural language reasoning with programmatic tool usage for solving problems, DeepSeekMath-Instruct 7B attains nearly 60\% accuracy on MATH, outperforming all previous open-source models. Across additional benchmarks, our model demonstrates competitive performance relative to DeepSeek-LLM-Chat 67B—the former leading model which possesses ten times as many parameters.  
\section{Reinforcement Learning}

\subsection{Group Relative Policy Optimization} Recent studies have demonstrated that reinforcement learning (RL) can significantly enhance the mathematical reasoning skills of LLMs following the Supervised Fine-Tuning (SFT) phase \citep{wang2023math,wizardmath}. Here, we present Group Relative Policy Optimization (GRPO), our robust and scalable RL technique designed for this purpose.

\subsubsection{From PPO to GRPO} Proximal Policy Optimization (PPO) \citep{schulman2017proximal} represents a widely adopted actor-critic reinforcement learning framework, particularly prevalent in the RL-based fine-tuning of large language models (LLMs) \citep{ouyang2022training}. This method seeks to enhance LLMs by maximizing the following surrogate loss:

\begin{equation}
\footnotesize
    \mathcal{J}_{PPO}(\theta) = \mathbb{E}{[q \sim P(Q), o \sim \pi_{\theta_{old}}(O|q)]} \frac{1}{|o|} \sum_{t=1}^{|o|} \min \left[ \frac{\pi_\theta(o_{t} | q, o_{<t})}{\pi_{\theta_{old}}(o_{t} | q, o_{<t})} A_{t}, \text{clip} \left( \frac{\pi_\theta(o_{t} | q, o_{<t})}{\pi_{\theta_{old}}(o_{t} | q, o_{<t})}, 1 - \varepsilon, 1 + \varepsilon \right)  A_{t} \right] ,
\end{equation}

In this context, $\pi_{\theta}$ represents the current policy, while $\pi_{\theta_{old}}$ denotes the previous version of the policy. The variables $q$ and $o$ correspond to questions and responses, respectively, drawn from the dataset of questions and outputs generated by the earlier policy $\pi_{\theta_{old}}$. The scalar $\varepsilon$ serves as a clipping hyper-parameter, an essential element in PPO that contributes to stable learning. The term $A_t$ refers to the advantage, calculated with Generalized Advantage Estimation (GAE) \citep{gae} using the sequence of rewards $\{r_{\ge t}\}$ together with an estimated value function $V_{\psi}$. Consequently, PPO requires the simultaneous training of a value function in addition to the policy network. To address the problem of reward model over-optimization, it is common practice to incorporate a token-wise KL divergence penalty—calculated using a reference model—into the reward for every generated token \citep{ouyang2022training}, that is,

\begin{equation}
    r_{t} = r_\varphi(q, o_{\le t}) - \beta \log\frac{\pi_{\theta}(o_{t}|q, o_{<t})}{\pi_{ref}(o_{t}|q, o_{<t})},
\label{eq:PPO-reward}
\end{equation}

where $r_\varphi$  is the reward model, $\pi_{ref}$ is the reference model, which is usually the initial SFT model, and $\beta$  is the coefficient of the KL penalty.

In PPO, the value function typically constitutes a model similar in size to the policy network, resulting in significant demands on memory and computation. Furthermore, throughout RL training, the value function serves as a baseline to compute the advantage for the purpose of reducing variance. In the specific scenario of LLMs, however, reward models frequently assign scores only to the final token, presenting a challenge for learning a value function that remains precise across all tokens. To circumvent these issues, as depicted in Figure \ref{fig:grpo}, we introduce Group Relative Policy Optimization (GRPO). GRPO eliminates the requirement for a dedicated value function model, in contrast to PPO, by leveraging the mean reward from a set of outputs generated in response to the same prompt as the baseline. Concretely, for a given question $q$, GRPO generates a collection of responses $\{o_1, o_2, \cdots, o_G\}$ from the previous policy $\pi_{\theta_{old}}$, and subsequently optimizes the policy network according to the objective function below:

\begin{equation}
\footnotesize
\begin{split}
    \mathcal{J}_{GRPO}(\theta) &= \mathbb{E}{[q \sim P(Q), \{o_i\}_{i=1}^G \sim \pi_{\theta_{old}}(O|q)]}  \\
    & \frac{1}{G} \sum_{i=1}^G \frac{1}{|o_i|} \sum_{t=1}^{|o_i|} \left\{ \min \left[ \frac{\pi_\theta(o_{i,t} | q, o_{i,<t})}{\pi_{\theta_{old}}(o_{i,t} | q, o_{i,<t})} \hat{A}_{i,t},\; \text{clip} \left( \frac{\pi_\theta(o_{i,t} | q, o_{i,<t})}{\pi_{\theta_{old}}(o_{i,t} | q, o_{i,<t})}, 1 - \varepsilon, 1 + \varepsilon \right) \hat{A}_{i,t} \right] - \beta \mathbb{D}_{KL}\big[\pi_{\theta} || \pi_{ref}\big] \right\} ,
\end{split}
\label{eq:GRPO-obj}
\end{equation}

Here, $\varepsilon$ and $\beta$ denote hyper-parameters, while $\hat{A}_{i,t}$ represents the advantage that is computed solely from the relative rewards of outputs within each respective group—a procedure elaborated upon in the subsequent sections. GRPO’s mechanism of calculating advantages based on group-relative measures is particularly suited to the inherent structure of reward models, which are generally trained using datasets built from pairwise output comparisons corresponding to the same prompt. Additionally, rather than introducing a KL penalty into the reward itself, GRPO imposes regularization by explicitly incorporating the KL divergence between the trained policy and the reference policy directly into the loss function. This circumvents the need to modify the computation of $\hat{A}_{i,t}$. Moreover, unlike the KL penalty employed in (\ref{eq:PPO-reward}), our approach utilizes the following unbiased estimator for the KL divergence \citep{kl_approx}:

\begin{equation}
\small
    \mathbb{D}_{KL}\left[\pi_{\theta} || \pi_{ref}\right] = \frac{\pi_{ref}(o_{i,t}|q,o_{i,<t})}{\pi_{\theta}(o_{i,t}|q,o_{i,<t})}- \log\frac{\pi_{ref}(o_{i,t}|q,o_{i,<t})}{\pi_{\theta}(o_{i,t}|q,o_{i,<t})} - 1,
\end{equation}

which is guaranteed to be positive.

\begin{algorithm}[t]
  \small
  \caption{Iterative Group Relative Policy Optimization}
  \textbf{Input} starting policy model $\pi_{\theta_{\text{init}}}$; reward estimators $r_\varphi$; collection of task prompts $\mathcal{D}$; 
  hyperparameters $\varepsilon$, $\beta$, $\mu$
  \begin{algorithmic}[1]
    \State Initialize policy model: $\pi_\theta \gets \pi_{\theta_{\text{init}}}$
    \For{iteration = 1, \dots, I}
      \State Set the reference model: $\pi_{ref} \gets \pi_{\theta}$
      \For{step = 1, \dots, M}
        \State Draw a minibatch $\mathcal{D}_b$ from $\mathcal{D}$
        \State Set $\pi_{\theta_{old}} \gets \pi_{\theta}$ 
        \State For each prompt $q$ in $\mathcal{D}_b$, generate $G$ samples $\{o_i\}_{i=1}^G$ from $\pi_{\theta_{old}}(\cdot|q)$
        \State Calculate the rewards $\{r_i\}_{i=1}^G$ for every output $o_i$ using $r_{\varphi}$
        \State For every output $o_i$, estimate the group-relative advantage $\hat{A}_{i,t}$ at each token position $t$
        \For{GRPO step = 1, \dots, $\mu$}
          \State Adjust $\pi_{\theta}$ by optimizing the GRPO objective (Equation \ref{eq:GC-GRPO})
        \EndFor
      \EndFor
      \State Refine $r_\varphi$ with continual learning based on a replay buffer.
    \EndFor
  \end{algorithmic}
  \textbf{Output} $\pi_\theta$
  \label{alg:iter-grpo}
\end{algorithm}

\subsubsection{Outcome Supervision RL with GRPO} In this approach, for each question $q$, a set of outputs $\{o_1, o_2, \cdots, o_G\}$ is generated by sampling from the old policy model $\pi_{\theta_{old}}$. These outputs are evaluated using a reward model, resulting in a set of $G$ rewards $\mathbf{r} = \{r_1, r_2, \cdots, r_G\}$. The rewards are then standardized by subtracting their mean and dividing by their standard deviation within the group. Under outcome supervision, the normalized reward is assigned to the end of each generated output $o_i$, and this normalized score serves as the advantage $\hat{A}_{i, t}$ for every token in the respective output. Specifically, $\hat{A}_{i, t} = \widetilde{r}_i = \frac{r_i- {\rm mean}(\mathbf{r})}{{\rm std}(\mathbf{r})}$. The policy is then updated through maximization of the objective specified in equation (\ref{eq:GRPO-obj}).

\subsubsection{Process Supervision RL with GRPO} Outcome supervision restricts feedback to a terminal reward after the model produces an output, which may be inadequate or inefficient for guiding policy learning in intricate mathematical reasoning. In line with \cite{wang2023math}, we investigate process supervision, which instead delivers a reward following each intermediate reasoning step. Given a question $q$ and $G$ sampled outputs $\{o_1, o_2, \cdots, o_G\}$, we employ a process reward model to evaluate every step within the outputs, generating step-level rewards: $\mathbf{R} = \{ \{r_1^{index(1)},\cdots,r_1^{index(K_1)}\}, \cdots,  \{r_G^{index(1)},\cdots,r_G^{index(K_G)}\} \}$, where $index(j)$ denotes the position of the end token in the $j$-th step and $K_i$ indicates the number of reasoning steps for output $i$. The rewards are standardized using their mean and standard deviation: $\widetilde{r}_i^{index(j)} = \frac{r_i^{index(j)} - {\rm mean(\mathbf{R})}}{{\rm std(\mathbf{R})}}$.

Process supervision then estimates the advantage for each token as the sum of subsequent normalized step rewards, specifically $\hat{A}_{i, t} = \sum_{index(j) \ge t} \widetilde{r}_i^{index(j)}$. Finally, the policy is updated by maximizing the objective introduced in equation (\ref{eq:GRPO-obj}).

\subsubsection{Iterative RL with GRPO} During reinforcement learning, the existing reward model may eventually fail to provide adequate supervision for the evolving policy model. To address this, we investigate an iterative approach for RL using GRPO. As described in Algorithm \ref{alg:iter-grpo}, this iterative procedure involves generating fresh training datasets for the reward model by sampling from the latest policy model. The reward model is incrementally updated, leveraging a replay strategy that maintains 10\% of the prior data for continual learning. Subsequently, the current policy model is used as the reference, and further training is performed on the policy model using the updated reward model.

\subsection{Training and Evaluating DeepSeekMath-RL}

Our reinforcement learning experiments utilize DeepSeekMath-Instruct 7B. The RL training data comprises chain-of-thought-style questions drawn from GSM8K and MATH subsets within the SFT dataset, totaling approximately 144K questions. To specifically study the effect of RL on benchmarks without direct training exposure during the RL phase, we omit all other SFT question types. The procedure for constructing the reward model training dataset follows the approach outlined in \citep{wang2023math}. We initialize the reward model using DeepSeekMath-Base 7B and adopt a learning rate of $2 \times 10^{-5}$. For the GRPO method, the policy model's learning rate is set to $1 \times 10^{-6}$, and a KL penalty coefficient of 0.04 is applied. We generate $64$ response samples per question. Both the maximum sequence length and the training batch size are configured to 1024. After each exploration iteration, the policy model undergoes only one update step. DeepSeekMath-RL 7B is evaluated on the same set of benchmarks as DeepSeekMath-Instruct 7B. Within this setup, GSM8K and MATH (with chain-of-thought annotations) are categorized as in-domain, while the rest of the benchmarks are treated as out-of-domain for DeepSeekMath-RL 7B.

Table \ref{tab:sft_rl_math} presents a comparative analysis of open- and closed-source models employing both chain-of-thought and tool-augmented reasoning across English and Chinese evaluation datasets. Our findings reveal the following: 1) DeepSeekMath-RL 7B achieves accuracy rates of 88.2\% on GSM8K and 51.7\% on MATH under chain-of-thought prompting, outperforming every open-source model within the 7B to 70B parameter range and exceeding the results of most closed-source systems as well. 2) Notably, DeepSeekMath-RL 7B undergoes training exclusively using chain-of-thought-formatted instruction data derived from GSM8K and MATH, initialized from DeepSeekMath-Instruct 7B. Despite the limited variety of its training data, DeepSeekMath-RL 7B consistently yields superior outcomes on all evaluated benchmarks compared to DeepSeekMath-Instruct 7B, thereby illustrating the advantages conferred by reinforcement learning.

\section{Discussion}
In this section, we will share our findings in pre-training and RL experiments. 

\subsection{Lessons Learnt in Pre-Training}

To begin, we describe our observations during the pre-training stage. Unless noted otherwise, the training configurations specified in Section \ref{sec:quality-policy} are consistently used. Additionally, throughout this section, references to the DeepSeekMath Corpus pertain to an 89B-token dataset assembled during the second round of data collection.

\subsubsection{Code Training Benefits Mathematical Reasoning}

An often-cited but unproven hypothesis posits that training on code enhances reasoning capabilities. In this work, we seek to partially address this claim within the context of mathematics, finding that exposure to code during training leads to improved mathematical reasoning in models, regardless of whether they employ external tools.

To study how code training affects mathematical reasoning, we experimented with the following two-stage training and one-stage training settings:

\textbf{Two-Stage Training}

\begin{itemize}[topsep=0pt]
    \item \textbf{Code Training for 400B Tokens $\rightarrow$ Math Training for 150B Tokens}: Our approach involves first training DeepSeek-LLM 1.3B on 400B tokens comprised of code, followed by an additional 150B tokens focused on mathematics;
    \item \textbf{General Training for 400B Tokens $\rightarrow$ Math Training for 150B Tokens}: As a comparative baseline, we conduct a parallel experiment where, in the initial phase, general-purpose tokens—drawn from an extensive, general corpus assembled by DeepSeek-AI—replace code tokens. This allows us to assess the relative benefits that code tokens may offer over general-domain tokens in boosting the model’s mathematical reasoning capabilities.
\end{itemize}

\textbf{One-Stage Training}

\begin{itemize}[topsep=0pt]
    \item \textbf{Mathematical Pretraining on 150B Tokens}: The DeepSeek-LLM 1.3B model undergoes training using 150B tokens from mathematical datasets.
    \item \textbf{Joint Training on a Blend of 400B Code Tokens and 150B Math Tokens}: Sequential math training after code finetuning diminishes performance on coding tasks. To address this, we examine if simultaneous exposure to both code and math tokens during a unified training phase can foster better mathematical reasoning capabilities while also mitigating catastrophic forgetting issues.
\end{itemize}

\paragraph{Results}
Table \ref{tab:code-math} and Table \ref{tab:code-math-general-eval} demonstrate the downstream performance under different training settings.

Incorporating code training facilitates program-aided mathematical reasoning in both one-stage and two-stage training regimes. Table \ref{tab:code-math} demonstrates that, within the two-stage training approach, solely training on code substantially improves performance on GSM8K and MATH problem-solving tasks executed in Python. Introducing a dedicated math training phase further advances these capabilities. Notably, when employing the one-stage training paradigm, integrating code and math tokens within the same phase effectively addresses the issue of catastrophic forgetting seen in two-stage setups, while also enhancing both coding skills (as shown in Table \ref{tab:code-math-general-eval}) and program-aided mathematical reasoning (see Table \ref{tab:code-math}).

Training on code additionally enhances mathematical reasoning skills even in the absence of tool assistance. Within the two-stage training paradigm, introducing code training at the outset already confers moderate gains. This phase further facilitates more efficient learning during the following math training stage, culminating in superior overall results. Conversely, merging code and math tokens during single-stage training appears to undermine mathematical reasoning when tools are not used. A possible explanation is that DeepSeek-LLM 1.3B, given its comparatively small model size, is unable to sufficiently integrate both code and mathematical content at the same time.

\subsubsection{ArXiv Papers Seem Ineffective in Improving Mathematical Reasoning}

ArXiv papers frequently serve as part of the mathematical pre-training datasets \citep{minerva,gpt-f,llemma,mathpile}. Despite this prevalence, there remains a paucity of thorough studies examining their influence on mathematical reasoning abilities. Somewhat unexpectedly, our empirical results indicate that arXiv papers do not substantially enhance mathematical reasoning skills. Our analysis spans a range of model sizes, such as DeepSeek-LLM 1.3B and DeepSeek-Coder-Base-v1.5 7B \citep{deepseek-coder}, utilizing arXiv datasets processed through a variety of distinct pipelines:

\begin{itemize}[topsep=0pt]
    \item \textbf{MathPile} \citep{mathpile}: a dataset consisting of 8.9B tokens, curated using specific cleaning and filtering strategies, where more than 85\% of the content originates from scientific arXiv manuscripts;
    \item \textbf{ArXiv-RedPajama} \citep{redpajama}: a collection comprising all arXiv LaTeX source files, processed to exclude preambles, comments, macros, and bibliographic entries, resulting in a total of 28.0B tokens.
\end{itemize}

Within our experimental framework, we independently fine-tune DeepSeek-LLM 1.3B on 150B tokens and DeepSeek-Coder-Base-v1.5 7B on 40B tokens for each respective arXiv dataset. Our observations indicate that exposure to arXiv papers alone does not enhance models' mathematical reasoning capabilities. When the models are trained exclusively on arXiv material, their performance remains largely unchanged or even declines across an array of mathematical evaluation benchmarks of varying levels of difficulty, as considered in this analysis. These assessments cover quantitative reasoning datasets such as GSM8K and MATH (Table \ref{tab:arxiv-cot}), multiple-choice formats exemplified by MMLU-STEM (Table \ref{tab:arxiv-cot}), as well as formal mathematics benchmarks like miniF2F (Table \ref{tab:arxiv_atp}).

However, this conclusion has its limitations and should be taken with a grain of salt. We have not yet studied:

\begin{itemize}[topsep=0pt]
    \item How arXiv tokens influence other mathematical applications beyond those analyzed here, for example, the informalization of theorems, which entails transforming formal statements or proofs into more informal language;
    \item The interaction between arXiv tokens and additional data modalities;
    \item If larger model architectures would reveal further advantages stemming from the inclusion of arXiv papers.
\end{itemize}

Thus, further exploration is required, which we leave for future studies.

\subsection{Insights of Reinforcement Learning}
\subsubsection{Towards to a Unified Paradigm} Here, we introduce a comprehensive framework to interpret a variety of training techniques—including SFT, RFT, DPO, PPO, and GRPO—by casting them within a single unified analysis. We then empirically investigate components of this unified structure. In most cases, the gradient of any specific training approach with respect to its parameter $\theta$ can be expressed as:

\begin{equation}
    \nabla_{\theta}\mathcal{J}_{\textcolor{red}{\mathcal{A}}}(\theta) = \mathbb{E}[\underbrace{(q,o) \sim \textcolor{red}{\mathcal{D}}}_{Data \ Source}]\left( \frac{1}{|o|} \sum_{t=1}^{|o|}  \underbrace{GC_{{\mathcal{A}}}(q, o, t, \textcolor{red}{\pi_{{rf}}})}_{Gradient \ Coefficient}  \nabla_{\theta}\log \pi_{\theta}(o_t | q, o_{<t})\right).
\label{eq:objective}
\end{equation}

Three principal elements are involved: 1) \textit{Data Source $\mathcal{D}$}, responsible for specifying the dataset used in training; 2) \textit{Reward Function $\pi_{{rf}}$}, which supplies the reward signal employed during training; and 3) \textit{Algorithm $\mathcal{A}$}, which leverages the training data along with the reward signal to compute the gradient coefficient $GC$, governing the extent of penalization or encouragement applied to the data. Utilizing this unified framework, we examine a number of representative techniques:

\begin{itemize}
    \item \textbf{Supervised Fine-tuning (SFT)}: SFT involves adapting a pretrained model through additional training on a curated dataset annotated by humans.
    \item \textbf{Rejection Sampling Fine-tuning (RFT)}: RFT provides another layer of fine-tuning on top of SFT by selecting and utilizing only the outputs generated by the SFT model that meet correctness criteria with respect to SFT question prompts.
    \item \textbf{Direct Preference Optimization (DPO)}: DPO enhances the SFT model by training it further on modified outputs produced by the SFT model, using a pairwise loss function specific to DPO.
    \item \textbf{Online Rejection Sampling Fine-tuning (Online RFT)}: In contrast to standard RFT, Online RFT initializes the policy from the SFT model but subsequently fine-tunes it using outputs generated from the live policy, rather than precomputed samples.
    \item \textbf{PPO/GRPO}: PPO/GRPO starts from an SFT-initialized policy model and applies reinforcement by leveraging outputs sampled dynamically from the current version of the policy model.
\end{itemize}

We summarize the components of these methods in Table \ref{tab:data_policy}.
Please refer to Appendix \ref{app:analysis-rl} for a more detailed derivation process.

\paragraph{Observation about Data Source} We categorize data sources into two main types: online sampling and offline sampling. Online sampling refers to obtaining training data from the exploration outcomes generated by the policy model as it is being trained. In contrast, offline sampling involves collecting training data from the sampling outputs of the pre-trained SFT model. Methods such as RFT and DPO employ the offline approach, whereas Online RFT and GRPO utilize the online approach.

Figure \ref{fig:rl-analysis} illustrates that Online RFT achieves markedly better results than RFT on two evaluation benchmarks. Initially, both methods perform similarly, but as training progresses, Online RFT attains a clear lead, highlighting the benefits of online training. This outcome is expected: during early training, the actor and SFT model remain quite similar, and the differences in their sampled data are minimal. In contrast, as training advances, divergences between the actor and the SFT model increase, making real-time data sampling considerably more beneficial.

\paragraph{Observation about Gradient Coefficient} The method utilizes the reward signal to compute the gradient coefficient, which in turn guides parameter updates of the model. In our experimental setup, we categorize the reward function into two types: ‘Rule’ and ‘Model.’ The ‘Rule’ category involves evaluating responses strictly on the basis of answer accuracy, whereas the ‘Model’ category makes use of a trained reward model to assign a score to each generated response. Notably, the training set for this reward model is constructed using outcomes derived from the rule-based evaluation. As presented in Equations \ref{eq:GC-RFT} and \ref{eq:GC-GRPO}, there is a significant distinction between GRPO and Online RFT methods: GRPO dynamically scales the gradient coefficient according to the reward output by the model, enabling nuanced differentiation in the degree of reinforcement or punishment based on the score magnitude. On the other hand, Online RFT does not incorporate this capability; it does not apply penalties to wrong responses and provides uniform reinforcement for all responses deemed correct, independent of their reward values.

Figure \ref{fig:rl-analysis} illustrates that GRPO outperforms online RFT, underscoring the effectiveness of modifying the coefficients for positive and negative gradients. Moreover, the results indicate that GRPO+PS achieves better outcomes than GRPO+OS, emphasizing the advantage of employing more granular, step-sensitive gradient coefficients. Additionally, we investigate the iterative RL approach by conducting two iterations in our experimental setup. The outcomes, depicted in Figure \ref{fig:iter-rl}, reveal that iterative RL leads to marked performance gains, with the most substantial improvement observed during the initial iteration.

\subsubsection{Why RL Works?} In this study, we apply reinforcement learning on a selected subset of instruction tuning data and observe that it yields notable gains over models trained only with instruction tuning. To provide further insight into the effectiveness of reinforcement learning, we assess both Pass@K and Maj@K accuracies for the Instruct and RL models using two evaluation benchmarks. Figure \ref{fig:maj-pass} illustrates that RL leads to improvements in Maj@K accuracy but does not significantly affect Pass@K. These outcomes suggest that RL bolsters the model’s performance primarily by making the output distribution more stable—specifically, \textbf{the benefits seem to result from increasing the frequency of correct answers within the TopK outputs, rather than enhancing the model’s intrinsic abilities.} In parallel, \citep{wang2023making} point out a \textbf{misalignment problem} in SFT models when tackling reasoning tasks, highlighting that their reasoning capabilities can be advanced through the application of preference alignment techniques \citep{yuan2023rrhf,song2023preference,wang2023making}.

\subsubsection{How to Achieve More Effective RL?}
\label{sec:effec_rl}
Our findings show that RL yields strong performance on mathematical reasoning tasks. Additionally, we introduce a unified framework that encapsulates various notable training strategies. In this framework, each approach is interpreted as a version of RL, whether in its pure form or in a simplified variant. As highlighted in Equation \ref{eq:objective}, the framework centers around three fundamental elements: Data Source, Algorithm, and Reward Function. We discuss possible avenues for advancing each of these three aspects.

\paragraph{Data Source} The origin of data constitutes the foundational input for all training approaches. Within the RL framework, the term “data source” is used to denote unlabeled questions accompanied by responses generated through sampling from the policy model. In this study, we restrict ourselves to utilizing questions provided during instruction tuning and apply a basic nucleus sampling for generating outputs. We hypothesize that this limitation may explain why our RL pipeline yields improvements solely in Maj@K scores. Moving forward, we intend to investigate the performance of our RL framework with out-of-distribution prompts, and to integrate \textbf{advanced sampling (decoding) strategies}—such as those employing tree-search methods \citep{yao2023tree}. Furthermore, the adoption of \textbf{efficient inference techniques} \citep{xia-etal-2023-speculative,leviathan2023fast,kwon2023efficient,xia2024unlocking}, which have a significant impact on the exploration capability of policy models, will also be an area of emphasis.

\paragraph{Algorithms} The role of algorithms is to leverage both data and reward signals in order to compute gradient coefficients, which are then used to adjust model parameters. As outlined in Equation \ref{eq:objective}, current methodologies generally place complete \textbf{TRUST} in the feedback provided by the reward function, modulating the conditional probability of specific tokens accordingly. Yet, this assumption that the reward signals are always trustworthy does not hold, particularly in highly intricate tasks. A case in point is the PRM800K dataset \citep{lightman2023let}, where despite extensive annotation by skilled professionals, about 20\% of the labels are still incorrect\footnote{\url{https://github.com/openai/prm800k/issues/12\#issuecomment-1728491852}}. Thus, it becomes essential to investigate reinforcement learning techniques that can maintain performance in the presence of noisy or unreliable reward signals. Our perspective is that \textbf{WEAK-TO-STRONG} alignment strategies \citep{burns2023weak} have the potential to transform learning algorithms in a profound way.

\paragraph{Reward Function} The reward function constitutes the primary feedback mechanism during training. In reinforcement learning, it is commonly instantiated as a neural reward model. We identify three critical research directions pertaining to reward models: 1) \textbf{Improving the generalization capacity of reward models.} Reward models need to generalize robustly to out-of-distribution prompts and complex generated outputs; a failure to do so risks that reinforcement learning only maintains LLM distributional behaviors without enhancing their underlying abilities; 2) \textbf{Capturing and representing the uncertainty inherent in the reward model.} Accounting for such uncertainty could serve as a valuable connector between imperfect reward models and methodologies that transfer from weaker to stronger learners; 3) \textbf{Developing scalable, high-fidelity process reward models} capable of delivering granular feedback during reasoning steps \citep{lightman2023let,wang2023math}.

\section{Conclusion, Limitation, and Future Work}

In this work, we introduce DeepSeekMath, a model that establishes new performance standards among open-source approaches on the competition-level MATH benchmark and achieves results close to those of proprietary systems. DeepSeekMath~is based on DeepSeek-Coder-v1.5 7B and is further trained over 500B tokens, including a substantial portion—120B tokens—of mathematical content extracted from Common Crawl. Through an extensive ablation analysis, we demonstrate that web pages constitute a promising resource for acquiring high-quality mathematical data, whereas data from arXiv appears less advantageous than initially anticipated. Additionally, we propose Group Relative Policy Optimization (GRPO), an adaptation of Proximal Policy Optimization (PPO), which significantly enhances mathematical reasoning performance while reducing memory requirements. Our experimental findings indicate that GRPO yields benefits even when DeepSeekMath-Instruct 7B achieves strong results on established benchmarks. Moreover, we outline a consolidated framework that contextualizes various methodological advancements and discuss several avenues for advancing reinforcement learning effectiveness in the domain.

While DeepSeekMath~demonstrates strong results on quantitative reasoning tasks, its performance on geometry and theorem-proving lags behind that of proprietary models. For example, during our preliminary testing, the model struggled with problems involving triangles and ellipses, suggesting possible pre-training and fine-tuning data selection biases. Additionally, the model's comparatively smaller scale limits its few-shot learning ability; unlike GPT-4, which benefits from few-shot examples, DeepSeekMath~displays little difference between its zero-shot and few-shot outcomes. Looking ahead, we intend to refine our data selection pipeline to assemble a higher-quality pre-training dataset. Moreover, we plan to investigate promising approaches (Section \ref{sec:effec_rl}) for advancing reinforcement learning in large language models.

\end{document}